% Part III -- Reusable and Reliable Programs
% Chapter 9: Copying Data Safely

\h{Copying Data Safely}

Two variables can refer to the same list or dictionary. When that happens, a
change made through one name appears through the other. This behavior is useful
when sharing is intentional and surprising when an independent copy was
expected.

\begin{NoteBox}
\textbf{Prerequisites.} You should understand variables, mutable collections,
nested lists and dictionaries, loops, and functions.
\end{NoteBox}

\begin{tcolorbox}[title=Learning outcomes,colframe=orange,colback=white,arc=2mm]
After completing this chapter, you should be able to:
\begin{itemize}
    \item distinguish equality from object identity;
    \item explain why assignment does not copy a mutable collection;
    \item make a shallow copy of a flat list or dictionary;
    \item identify shared inner objects in nested data;
    \item use \c{copy.deepcopy()} when full independence is required; and
    \item prevent unintended mutation inside functions.
\end{itemize}
\end{tcolorbox}

\hh{Equality and identity answer different questions}

The \c{==} operator asks whether two values are equal. The \c{is} operator asks
whether two names refer to the same object.

\code{c09_equality_identity}{python}{
    first = ["Lobby", "Office"]
    second = ["Lobby", "Office"]

    print(first == second)
    print(first is second)
}{Compare equal values stored in separate lists}

The lists contain equal values, so \c{==} returns \c{True}. They are separate
list objects, so \c{is} returns \c{False}.

Use \c{==} for normal value comparisons. Use \c{is} mainly for identity checks,
especially \c{value is None}.

\hh{Assignment creates another name}

Assigning a list to another variable does not create a new list.

\code{c09_assignment_alias}{python}{
    original_rooms = ["Lobby", "Office"]
    shared_rooms = original_rooms

    shared_rooms.append("Cafe")

    print(original_rooms)
    print(shared_rooms)
    print(original_rooms is shared_rooms)
}{Two names refer to one list}

Both names point to the same mutable object. Appending through either name
changes that one list.

\begin{NoteBox}
\textbf{Stop and predict.} If \c{original_rooms.append("Gallery")} runs next,
what will \c{shared_rooms} contain?
\end{NoteBox}

\hh{Immutable values behave differently}

Numbers and strings cannot be changed in place. Reassignment connects a name to
a different value; it does not change the earlier value.

\code{c09_immutable_reassignment}{python}{
    first_count = 3
    second_count = first_count

    second_count = 4

    print(first_count)
    print(second_count)
}{Reassign one name bound to an integer}

The output is 3 and 4. No integer object was mutated.

\hh{Make a shallow list copy}

For a flat list, \c{.copy()}, a full slice, and \c{list()} each create a new
outer list.

\code{c09_flat_list_copy}{python}{
    original_rooms = ["Lobby", "Office"]
    copied_rooms = original_rooms.copy()

    copied_rooms.append("Cafe")

    print(original_rooms)
    print(copied_rooms)
    print(original_rooms is copied_rooms)
}{Copy a flat list}

The copied list is independent because the contained strings are immutable.

\hh{Make shallow dictionary and set copies}

The \c{.copy()} method also creates new dictionary and set containers.

\code{c09_dict_set_copy}{python}{
    original_room = {"name": "Lobby", "area_m2": 24.0}
    copied_room = original_room.copy()
    copied_room["area_m2"] = 25.0

    original_levels = {1, 2, 3}
    copied_levels = original_levels.copy()
    copied_levels.add(4)

    print(original_room)
    print(copied_room)
    print(original_levels)
    print(copied_levels)
}{Copy flat dictionaries and sets}

These examples are flat. Their values do not contain mutable collections.

\hh{A shallow copy shares nested objects}

A shallow copy creates a new outer container but reuses the objects inside it.

\code{c09_shallow_nested}{python}{
    original = {
        "name": "Lobby",
        "tags": ["public", "entry"]
    }

    shallow = original.copy()
    shallow["tags"].append("display")

    print(original["tags"])
    print(shallow["tags"])
    print(original["tags"] is shallow["tags"])
}{A shallow dictionary copy shares its inner list}

The dictionaries are separate, but their \c{"tags"} values point to the same
list. That is the meaning of \emph{shallow}.

\hh{Copy only the nested part that will change}

Sometimes a targeted copy is clearer than copying the entire structure.

\code{c09_targeted_nested_copy}{python}{
    original = {
        "name": "Lobby",
        "tags": ["public", "entry"]
    }

    revised = original.copy()
    revised["tags"] = original["tags"].copy()
    revised["tags"].append("display")

    print(original["tags"])
    print(revised["tags"])
}{Copy the outer dictionary and one inner list}

This approach states exactly which nested value must become independent.

\hh{Use a deep copy for full nested independence}

The standard-library \c{copy} module can recursively duplicate nested mutable
objects.

\code{c09_deep_copy}{python}{
    import copy

    original = {
        "name": "Lobby",
        "tags": ["public", "entry"],
        "dimensions_m": {"width": 6.0, "length": 8.0}
    }

    independent = copy.deepcopy(original)
    independent["tags"].append("display")
    independent["dimensions_m"]["width"] = 7.0

    print(original)
    print(independent)
}{Deep-copy nested data before independent editing}

Deep copying is convenient, but it may copy more than necessary. Prefer a
targeted copy when the structure is large and only one part will change.

\hh{Avoid replicated inner lists}

Multiplying a list that contains another list repeats references to the same
inner object.

\code{c09_grid_alias_problem}{python}{
    grid = [[0, 0, 0]] * 3
    grid[0][0] = 1

    print(grid)
}{A replicated row is shared three times}

Create each row separately instead.

\code{c09_grid_independent_rows}{python}{
    grid = []

    for _ in range(3):
        grid.append([0, 0, 0])

    grid[0][0] = 1
    print(grid)
}{Build independent rows with an ordinary loop}

The underscore is a conventional loop variable when its value is not used.

\hh{Functions can mutate caller data}

A function receives a reference to the supplied list. A method such as
\c{append()} therefore changes the caller's list.

\code{c09_function_mutation}{python}{
    def add_room_in_place(room_names, new_name):
        room_names.append(new_name)

    rooms = ["Lobby", "Office"]
    add_room_in_place(rooms, "Cafe")

    print(rooms)
}{A function intentionally mutates its input list}

The function name includes \c{in_place} to make the mutation visible. When the
caller should keep its data, copy before changing it.

\code{c09_function_copy}{python}{
    def with_added_room(room_names, new_name):
        revised_names = room_names.copy()
        revised_names.append(new_name)
        return revised_names

    rooms = ["Lobby", "Office"]
    expanded_rooms = with_added_room(rooms, "Cafe")

    print(rooms)
    print(expanded_rooms)
}{Return a revised copy and preserve the input}

Neither approach is automatically correct. The important point is to make the
behavior intentional and clear.

\hh{Choose the smallest sufficient copy}

\begin{center}
\begin{tabular}{ll}
\textbf{Situation} & \textbf{Usual choice} \\
Two names should share one object & Assignment \\
Flat list, dictionary, or set & Shallow copy \\
Nested data; only one inner value changes & Targeted shallow copies \\
Many nested values must be independent & \c{copy.deepcopy()} \\
Immutable number or string & Usually no copy needed \\
\end{tabular}
\end{center}

Inspect the structure and planned changes before choosing.

\hh{Common mistakes}

\begin{itemize}
    \item \textbf{Using assignment when independence is expected}: both names share changes.
    \item \textbf{Assuming a shallow copy duplicates nested lists}: inner objects remain shared.
    \item \textbf{Using is for value equality}: use \c{==} for normal comparisons.
    \item \textbf{Deep-copying everything automatically}: it may be unnecessary and expensive.
    \item \textbf{Mutating a function argument without saying so}: callers may be surprised.
    \item \textbf{Replicating a nested list with multiplication}: rows may share one object.
\end{itemize}

\hh{Practice ladder}

\hhh{Level 0 -- Read and predict}

\code{c09_practice_level0}{python}{
    first = [1, 2]
    second = first
    second.append(3)

    print(first)
    print(first is second)
}{Predict shared-list output}

\hhh{Level 1 -- Modify}

Change the assignment to \c{second = first.copy()}. Predict both outputs again.

\hhh{Level 2 -- Complete}

Create an independent copy of a flat dictionary and change only the copy.

\code{c09_practice_level2}{python}{
    original = {"name": "Lobby", "area_m2": 24.0}
    revised = original.copy()
    revised["area_m2"] = 26.0

    print(original)
    print(revised)
}{Copy and revise a flat dictionary}

\hhh{Level 3 -- Apply}

Given a dictionary containing a room name and a list of tags, create a revised
version with one additional tag while preserving the original.

\textbf{Hints}
\begin{itemize}
    \item \textbf{Hint 1}: A dictionary copy alone is shallow.
    \item \textbf{Hint 2}: Copy the outer dictionary and the inner tag list.
    \item \textbf{Hint 3}: Append only after both copies exist.
\end{itemize}

\hhh{Level 4 -- Challenge}

Write \c{with_updated_width(room, width_m)}. Return an independent nested
dictionary with the revised width while preserving all original data.

\hh{Practice solutions}

\textbf{Level 0:} The outputs are \c{[1, 2, 3]} and \c{True}.

\textbf{Level 2:} Only \c{revised} contains the area \c{26.0}.

\textbf{Level 3}

\code{c09_solution_level3}{python}{
    original = {
        "name": "Lobby",
        "tags": ["public", "entry"]
    }

    revised = original.copy()
    revised["tags"] = original["tags"].copy()
    revised["tags"].append("display")

    print(original)
    print(revised)
}{One solution to Level 3}

\textbf{Level 4}

\code{c09_solution_level4}{python}{
    import copy

    def with_updated_width(room, width_m):
        revised = copy.deepcopy(room)
        revised["dimensions_m"]["width"] = width_m
        return revised

    original = {
        "name": "Lobby",
        "dimensions_m": {"width": 6.0, "length": 8.0}
    }

    revised = with_updated_width(original, 7.0)
    print(original)
    print(revised)
}{One solution to Level 4}

\hh{Chapter checkpoint}

\begin{enumerate}
    \item What question does \c{==} answer?
    \item What question does \c{is} answer?
    \item Does assigning a list to a new name copy the list?
    \item What does a shallow copy duplicate?
    \item What remains shared in a shallow copy of nested data?
    \item When is a targeted copy useful?
    \item What does \c{copy.deepcopy()} do?
    \item Can a function mutate a list supplied by its caller?
    \item Why is \c{[[0, 0]] * 3} risky?
\end{enumerate}

\hh{Checkpoint answers}

\begin{enumerate}
    \item Whether two values are equal.
    \item Whether two names refer to the same object.
    \item No. It creates another reference to the same list.
    \item The outer container.
    \item The inner mutable objects.
    \item When only selected nested parts need independence.
    \item It recursively copies nested objects.
    \item Yes, if it performs an in-place change.
    \item Each repeated row may refer to the same inner list.
\end{enumerate}

\begin{tcolorbox}[title=Chapter summary,colframe=orange,colback=white,arc=2mm]
You can distinguish equality from identity, recognize aliases, make shallow and
deep copies, copy selected nested values, and control mutation in functions.
The next chapter stores and retrieves data with files and paths.
\end{tcolorbox}

